
\documentclass{article}
\usepackage{colortbl}
\usepackage{makecell}
\usepackage{multirow}
\usepackage{supertabular}

\begin{document}

\newcounter{utterance}

\centering \large Interaction Transcript for game `dond', experiment `coop\_it', episode 14 with claude{-}sonnet{-}4{-}20250514{-}no{-}reasoning{-}t1.0.
\vspace{24pt}

{ \footnotesize  \setcounter{utterance}{1}
\setlength{\tabcolsep}{0pt}
\begin{supertabular}{c@{$\;$}|p{.15\linewidth}@{}p{.15\linewidth}p{.15\linewidth}p{.15\linewidth}p{.15\linewidth}p{.15\linewidth}}
    \# & $\;$A & \multicolumn{4}{c}{Game Master} & $\;\:$B\\
    \hline

    \theutterance \stepcounter{utterance}  
    & & \multicolumn{4}{p{0.6\linewidth}}{
        \cellcolor[rgb]{0.9,0.9,0.9}{
            \makecell[{{p{\linewidth}}}]{
                \texttt{\tiny{[P1$\langle$GM]}}
                \texttt{State giocando a un gioco di negoziazione cooperativa in cui dovete accordarvi su come dividere una serie di oggetti tra voi e un altro giocatore.} \\
\\ 
\texttt{Regole:} \\
\texttt{(a) A Lei e all'altro giocatore viene dato un insieme di oggetti. Ognuno di voi riceve anche una funzione di valore segreta, che rappresenta il valore di ciascun tipo di oggetto.} \\
\texttt{(b) Si scambiano messaggi con l'altro giocatore per concordare chi si aggiudica gli oggetti. Potete inviare un massimo di 5 messaggi ciascuno, oppure terminare in anticipo facendo una proposta segreta in qualsiasi momento.} \\
\texttt{(c) A ciascuno di voi viene chiesto di inviare una proposta segreta indicando gli oggetti che desiderate, formattata tra parentesi quadre come segue: "{[}Proposta: <numero> <nome oggetto>, <numero> <nome oggetto>, <...>{]}".} \\
\texttt{(d) Se le vostre proposte sono complementari, cioè ci sono abbastanza oggetti per soddisfare entrambe le proposte, a ciascun giocatore viene assegnato un punteggio basato sulla somma dei valori degli oggetti ricevuti. In caso contrario, entrambi ricevono zero punti.} \\
\texttt{(e) L'obiettivo è massimizzare la somma del proprio punteggio e di quello dell'altro giocatore.} \\
\\ 
\texttt{Cominciamo.} \\
\\ 
\texttt{L'insieme degli oggetti disponibili è:} \\
\\ 
\texttt{2 specchi, 1 mattone, 3 coltelli, 1 pettine} \\
\\ 
\texttt{La funzione valore segreta è:} \\
\\ 
\texttt{specchio: 1, mattone: 4, coltello: 1, pettine: 1} \\
\\ 
\texttt{IMPORTANTE: i vostri messaggi, a meno che non si tratti di una proposta segreta, vengono trasmessi direttamente all'altro giocatore, quindi non includete alcuna risposta alle regole o testo di annuncio del vostro messaggio. Per fare una proposta segreta, utilizzate il formato indicato. Non utilizzare delle parentesi quadre quando si comunica all'altro giocatore, altrimenti verrà interpretata come una proposta segreta.} \\
            }
        }
    }
    & \\ \\

    \theutterance \stepcounter{utterance}  
    & \multicolumn{4}{p{0.6\linewidth}}{
        \cellcolor[rgb]{0.9,0.9,0.9}{
            \makecell[{{p{\linewidth}}}]{
                \texttt{\tiny{[P1$\rangle$GM]}}
                \texttt{Ciao! Vediamo cosa abbiamo da dividere: 2 specchi, 1 mattone, 3 coltelli e 1 pettine. Per trovare la divisione migliore per entrambi, dovremmo capire quali oggetti sono più importanti per ciascuno di noi. } \\
\\ 
\texttt{Puoi dirmi quali oggetti preferisci o hanno più valore per te? Io ho delle preferenze specifiche ma vorrei sentire le tue prima di fare proposte.} \\
            }
        }
    }
    & & \\ \\

    \theutterance \stepcounter{utterance}  
    & & & \multicolumn{2}{p{0.3\linewidth}}{
        \cellcolor[rgb]{0.9,0.9,0.9}{
            \makecell[{{p{\linewidth}}}]{
                \texttt{\tiny{[GM$|$GM]}}
                \texttt{continue} \\
            }
        }
    }
    & & \\ \\

    \theutterance \stepcounter{utterance}  
    & & & \multicolumn{2}{p{0.3\linewidth}}{
        \cellcolor[rgb]{0.9,0.9,0.9}{
            \makecell[{{p{\linewidth}}}]{
                \texttt{\tiny{[GM$|$GM]}}
                \texttt{Ciao! Vediamo cosa abbiamo da dividere: 2 specchi, 1 mattone, 3 coltelli e 1 pettine. Per trovare la divisione migliore per entrambi, dovremmo capire quali oggetti sono più importanti per ciascuno di noi. } \\
\\ 
\texttt{Puoi dirmi quali oggetti preferisci o hanno più valore per te? Io ho delle preferenze specifiche ma vorrei sentire le tue prima di fare proposte.} \\
            }
        }
    }
    & & \\ \\

    \theutterance \stepcounter{utterance}  
    & & \multicolumn{4}{p{0.6\linewidth}}{
        \cellcolor[rgb]{0.9,0.9,0.9}{
            \makecell[{{p{\linewidth}}}]{
                \texttt{\tiny{[GM$\rangle$P2]}}
                \texttt{State giocando a un gioco di negoziazione cooperativa in cui dovete accordarvi su come dividere una serie di oggetti tra voi e un altro giocatore.} \\
\\ 
\texttt{Regole:} \\
\texttt{(a) A Lei e all'altro giocatore viene dato un insieme di oggetti. Ognuno di voi riceve anche una funzione di valore segreta, che rappresenta il valore di ciascun tipo di oggetto.} \\
\texttt{(b) Si scambiano messaggi con l'altro giocatore per concordare chi si aggiudica gli oggetti. Potete inviare un massimo di 5 messaggi ciascuno, oppure terminare in anticipo facendo una proposta segreta in qualsiasi momento.} \\
\texttt{(c) A ciascuno di voi viene chiesto di inviare una proposta segreta indicando gli oggetti che desiderate, formattata tra parentesi quadre come segue: "{[}Proposta: <numero> <nome oggetto>, <numero> <nome oggetto>, <...>{]}".} \\
\texttt{(d) Se le vostre proposte sono complementari, cioè ci sono abbastanza oggetti per soddisfare entrambe le proposte, a ciascun giocatore viene assegnato un punteggio basato sulla somma dei valori degli oggetti ricevuti. In caso contrario, entrambi ricevono zero punti.} \\
\texttt{(e) L'obiettivo è massimizzare la somma del proprio punteggio e di quello dell'altro giocatore.} \\
\\ 
\texttt{Cominciamo.} \\
\\ 
\texttt{L'insieme degli oggetti disponibili è:} \\
\\ 
\texttt{2 specchi, 1 mattone, 3 coltelli, 1 pettine} \\
\\ 
\texttt{La funzione valore segreta è:} \\
\\ 
\texttt{specchio: 2, mattone: 1, coltello: 0, pettine: 5} \\
\\ 
\texttt{IMPORTANTE: i vostri messaggi, a meno che non si tratti di una proposta segreta, vengono trasmessi direttamente all'altro giocatore, quindi non includete alcuna risposta alle regole o testo di annuncio del vostro messaggio. Per fare una proposta segreta, utilizzate il formato indicato. Non utilizzare delle parentesi quadre quando si comunica all'altro giocatore, altrimenti verrà interpretata come una proposta segreta.} \\
\\ 
\\ 
\texttt{Ciao! Vediamo cosa abbiamo da dividere: 2 specchi, 1 mattone, 3 coltelli e 1 pettine. Per trovare la divisione migliore per entrambi, dovremmo capire quali oggetti sono più importanti per ciascuno di noi. } \\
\\ 
\texttt{Puoi dirmi quali oggetti preferisci o hanno più valore per te? Io ho delle preferenze specifiche ma vorrei sentire le tue prima di fare proposte.} \\
            }
        }
    }
    & \\ \\

    \theutterance \stepcounter{utterance}  
    & & & \multicolumn{4}{p{0.6\linewidth}}{
        \cellcolor[rgb]{0.9,0.9,0.9}{
            \makecell[{{p{\linewidth}}}]{
                \texttt{\tiny{[GM$\langle$P2]}}
                \texttt{Ciao! Sono d'accordo che dovremmo capire le reciproche preferenze per trovare una buona divisione.} \\
\\ 
\texttt{Per me, in ordine di preferenza, direi che i mattoni sono molto utili, seguiti dagli specchi che hanno un buon valore pratico. I pettini li trovo meno interessanti, e i coltelli sono quelli che mi servono di meno.} \\
\\ 
\texttt{E tu? Quali oggetti ti interessano di più? Se riusciamo a capire dove le nostre preferenze si differenziano, possiamo probabilmente trovare una divisione che funzioni bene per entrambi.} \\
            }
        }
    }
     \\ \\

    \theutterance \stepcounter{utterance}  
    & & & \multicolumn{2}{p{0.3\linewidth}}{
        \cellcolor[rgb]{0.9,0.9,0.9}{
            \makecell[{{p{\linewidth}}}]{
                \texttt{\tiny{[GM$|$GM]}}
                \texttt{continue} \\
            }
        }
    }
    & & \\ \\

    \theutterance \stepcounter{utterance}  
    & & & \multicolumn{2}{p{0.3\linewidth}}{
        \cellcolor[rgb]{0.9,0.9,0.9}{
            \makecell[{{p{\linewidth}}}]{
                \texttt{\tiny{[GM$|$GM]}}
                \texttt{Ciao! Sono d'accordo che dovremmo capire le reciproche preferenze per trovare una buona divisione.} \\
\\ 
\texttt{Per me, in ordine di preferenza, direi che i mattoni sono molto utili, seguiti dagli specchi che hanno un buon valore pratico. I pettini li trovo meno interessanti, e i coltelli sono quelli che mi servono di meno.} \\
\\ 
\texttt{E tu? Quali oggetti ti interessano di più? Se riusciamo a capire dove le nostre preferenze si differenziano, possiamo probabilmente trovare una divisione che funzioni bene per entrambi.} \\
            }
        }
    }
    & & \\ \\

    \theutterance \stepcounter{utterance}  
    & & \multicolumn{4}{p{0.6\linewidth}}{
        \cellcolor[rgb]{0.9,0.9,0.9}{
            \makecell[{{p{\linewidth}}}]{
                \texttt{\tiny{[P1$\langle$GM]}}
                \texttt{Ciao! Sono d'accordo che dovremmo capire le reciproche preferenze per trovare una buona divisione.} \\
\\ 
\texttt{Per me, in ordine di preferenza, direi che i mattoni sono molto utili, seguiti dagli specchi che hanno un buon valore pratico. I pettini li trovo meno interessanti, e i coltelli sono quelli che mi servono di meno.} \\
\\ 
\texttt{E tu? Quali oggetti ti interessano di più? Se riusciamo a capire dove le nostre preferenze si differenziano, possiamo probabilmente trovare una divisione che funzioni bene per entrambi.} \\
            }
        }
    }
    & \\ \\

    \theutterance \stepcounter{utterance}  
    & \multicolumn{4}{p{0.6\linewidth}}{
        \cellcolor[rgb]{0.9,0.9,0.9}{
            \makecell[{{p{\linewidth}}}]{
                \texttt{\tiny{[P1$\rangle$GM]}}
                \texttt{Perfetto! Le nostre preferenze sembrano abbastanza complementari. Anch'io trovo il mattone molto utile, ma vedo che è importante anche per te. Per quanto riguarda gli altri oggetti, sono abbastanza flessibile.} \\
\\ 
\texttt{Dato che abbiamo entrambi interesse per il mattone ma ce n'è solo uno, forse possiamo bilanciare la divisione in questo modo: chi prende il mattone potrebbe lasciare all'altro una quota maggiore degli altri oggetti?} \\
\\ 
\texttt{Ad esempio, potrei prendere il mattone e forse uno specchio, mentre tu potresti avere l'altro specchio, tutti e 3 i coltelli (che dici ti servono meno ma comunque hanno un certo valore), e il pettine. } \\
\\ 
\texttt{Che ne pensi di questa divisione? Ti sembra equa o preferiresti aggiustare qualcosa?} \\
            }
        }
    }
    & & \\ \\

    \theutterance \stepcounter{utterance}  
    & & & \multicolumn{2}{p{0.3\linewidth}}{
        \cellcolor[rgb]{0.9,0.9,0.9}{
            \makecell[{{p{\linewidth}}}]{
                \texttt{\tiny{[GM$|$GM]}}
                \texttt{continue} \\
            }
        }
    }
    & & \\ \\

    \theutterance \stepcounter{utterance}  
    & & & \multicolumn{2}{p{0.3\linewidth}}{
        \cellcolor[rgb]{0.9,0.9,0.9}{
            \makecell[{{p{\linewidth}}}]{
                \texttt{\tiny{[GM$|$GM]}}
                \texttt{Perfetto! Le nostre preferenze sembrano abbastanza complementari. Anch'io trovo il mattone molto utile, ma vedo che è importante anche per te. Per quanto riguarda gli altri oggetti, sono abbastanza flessibile.} \\
\\ 
\texttt{Dato che abbiamo entrambi interesse per il mattone ma ce n'è solo uno, forse possiamo bilanciare la divisione in questo modo: chi prende il mattone potrebbe lasciare all'altro una quota maggiore degli altri oggetti?} \\
\\ 
\texttt{Ad esempio, potrei prendere il mattone e forse uno specchio, mentre tu potresti avere l'altro specchio, tutti e 3 i coltelli (che dici ti servono meno ma comunque hanno un certo valore), e il pettine. } \\
\\ 
\texttt{Che ne pensi di questa divisione? Ti sembra equa o preferiresti aggiustare qualcosa?} \\
            }
        }
    }
    & & \\ \\

    \theutterance \stepcounter{utterance}  
    & & \multicolumn{4}{p{0.6\linewidth}}{
        \cellcolor[rgb]{0.9,0.9,0.9}{
            \makecell[{{p{\linewidth}}}]{
                \texttt{\tiny{[GM$\rangle$P2]}}
                \texttt{Perfetto! Le nostre preferenze sembrano abbastanza complementari. Anch'io trovo il mattone molto utile, ma vedo che è importante anche per te. Per quanto riguarda gli altri oggetti, sono abbastanza flessibile.} \\
\\ 
\texttt{Dato che abbiamo entrambi interesse per il mattone ma ce n'è solo uno, forse possiamo bilanciare la divisione in questo modo: chi prende il mattone potrebbe lasciare all'altro una quota maggiore degli altri oggetti?} \\
\\ 
\texttt{Ad esempio, potrei prendere il mattone e forse uno specchio, mentre tu potresti avere l'altro specchio, tutti e 3 i coltelli (che dici ti servono meno ma comunque hanno un certo valore), e il pettine. } \\
\\ 
\texttt{Che ne pensi di questa divisione? Ti sembra equa o preferiresti aggiustare qualcosa?} \\
            }
        }
    }
    & \\ \\

    \theutterance \stepcounter{utterance}  
    & & & \multicolumn{4}{p{0.6\linewidth}}{
        \cellcolor[rgb]{0.9,0.9,0.9}{
            \makecell[{{p{\linewidth}}}]{
                \texttt{\tiny{[GM$\langle$P2]}}
                \texttt{Hmm, apprezzo la tua proposta, ma devo essere onesto {-} anche se ho detto che i coltelli mi servono meno, prenderne 3 insieme al pettine mi sembra un po' sbilanciato considerando che tu avresti sia il mattone che uno specchio.} \\
\\ 
\texttt{Che ne dici se invece facciamo così: tu prendi il mattone e il pettine (visto che hai detto di essere flessibile), e io prendo i 2 specchi e tutti i 3 coltelli? } \\
\\ 
\texttt{In questo modo tu avresti comunque il mattone che entrambi apprezziamo, più un oggetto aggiuntivo, mentre io avrei una quantità maggiore di oggetti ma di valore per me inferiore. Mi sembra più equilibrato. Che ne pensi?} \\
            }
        }
    }
     \\ \\

    \theutterance \stepcounter{utterance}  
    & & & \multicolumn{2}{p{0.3\linewidth}}{
        \cellcolor[rgb]{0.9,0.9,0.9}{
            \makecell[{{p{\linewidth}}}]{
                \texttt{\tiny{[GM$|$GM]}}
                \texttt{continue} \\
            }
        }
    }
    & & \\ \\

    \theutterance \stepcounter{utterance}  
    & & & \multicolumn{2}{p{0.3\linewidth}}{
        \cellcolor[rgb]{0.9,0.9,0.9}{
            \makecell[{{p{\linewidth}}}]{
                \texttt{\tiny{[GM$|$GM]}}
                \texttt{Hmm, apprezzo la tua proposta, ma devo essere onesto {-} anche se ho detto che i coltelli mi servono meno, prenderne 3 insieme al pettine mi sembra un po' sbilanciato considerando che tu avresti sia il mattone che uno specchio.} \\
\\ 
\texttt{Che ne dici se invece facciamo così: tu prendi il mattone e il pettine (visto che hai detto di essere flessibile), e io prendo i 2 specchi e tutti i 3 coltelli? } \\
\\ 
\texttt{In questo modo tu avresti comunque il mattone che entrambi apprezziamo, più un oggetto aggiuntivo, mentre io avrei una quantità maggiore di oggetti ma di valore per me inferiore. Mi sembra più equilibrato. Che ne pensi?} \\
            }
        }
    }
    & & \\ \\

    \theutterance \stepcounter{utterance}  
    & & \multicolumn{4}{p{0.6\linewidth}}{
        \cellcolor[rgb]{0.9,0.9,0.9}{
            \makecell[{{p{\linewidth}}}]{
                \texttt{\tiny{[P1$\langle$GM]}}
                \texttt{Hmm, apprezzo la tua proposta, ma devo essere onesto {-} anche se ho detto che i coltelli mi servono meno, prenderne 3 insieme al pettine mi sembra un po' sbilanciato considerando che tu avresti sia il mattone che uno specchio.} \\
\\ 
\texttt{Che ne dici se invece facciamo così: tu prendi il mattone e il pettine (visto che hai detto di essere flessibile), e io prendo i 2 specchi e tutti i 3 coltelli? } \\
\\ 
\texttt{In questo modo tu avresti comunque il mattone che entrambi apprezziamo, più un oggetto aggiuntivo, mentre io avrei una quantità maggiore di oggetti ma di valore per me inferiore. Mi sembra più equilibrato. Che ne pensi?} \\
            }
        }
    }
    & \\ \\

    \theutterance \stepcounter{utterance}  
    & \multicolumn{4}{p{0.6\linewidth}}{
        \cellcolor[rgb]{0.9,0.9,0.9}{
            \makecell[{{p{\linewidth}}}]{
                \texttt{\tiny{[P1$\rangle$GM]}}
                \texttt{Sì, questa divisione mi sembra molto più equilibrata! Hai ragione che la mia prima proposta era un po' sbilanciata.} \\
\\ 
\texttt{La tua proposta ha senso: tu avresti 2 specchi e 3 coltelli (5 oggetti in totale), mentre io avrei il mattone e il pettine (2 oggetti). Considerando che il mattone sembra essere l'oggetto di maggior valore per entrambi, è giusto che chi lo prende abbia meno oggetti totali.} \\
\\ 
\texttt{Sono d'accordo con questa divisione: tu prendi 2 specchi e 3 coltelli, io prendo 1 mattone e 1 pettine.} \\
\\ 
\texttt{Ti va bene se procediamo con questa proposta?} \\
            }
        }
    }
    & & \\ \\

    \theutterance \stepcounter{utterance}  
    & & & \multicolumn{2}{p{0.3\linewidth}}{
        \cellcolor[rgb]{0.9,0.9,0.9}{
            \makecell[{{p{\linewidth}}}]{
                \texttt{\tiny{[GM$|$GM]}}
                \texttt{continue} \\
            }
        }
    }
    & & \\ \\

    \theutterance \stepcounter{utterance}  
    & & & \multicolumn{2}{p{0.3\linewidth}}{
        \cellcolor[rgb]{0.9,0.9,0.9}{
            \makecell[{{p{\linewidth}}}]{
                \texttt{\tiny{[GM$|$GM]}}
                \texttt{Sì, questa divisione mi sembra molto più equilibrata! Hai ragione che la mia prima proposta era un po' sbilanciata.} \\
\\ 
\texttt{La tua proposta ha senso: tu avresti 2 specchi e 3 coltelli (5 oggetti in totale), mentre io avrei il mattone e il pettine (2 oggetti). Considerando che il mattone sembra essere l'oggetto di maggior valore per entrambi, è giusto che chi lo prende abbia meno oggetti totali.} \\
\\ 
\texttt{Sono d'accordo con questa divisione: tu prendi 2 specchi e 3 coltelli, io prendo 1 mattone e 1 pettine.} \\
\\ 
\texttt{Ti va bene se procediamo con questa proposta?} \\
            }
        }
    }
    & & \\ \\

    \theutterance \stepcounter{utterance}  
    & & \multicolumn{4}{p{0.6\linewidth}}{
        \cellcolor[rgb]{0.9,0.9,0.9}{
            \makecell[{{p{\linewidth}}}]{
                \texttt{\tiny{[GM$\rangle$P2]}}
                \texttt{Sì, questa divisione mi sembra molto più equilibrata! Hai ragione che la mia prima proposta era un po' sbilanciata.} \\
\\ 
\texttt{La tua proposta ha senso: tu avresti 2 specchi e 3 coltelli (5 oggetti in totale), mentre io avrei il mattone e il pettine (2 oggetti). Considerando che il mattone sembra essere l'oggetto di maggior valore per entrambi, è giusto che chi lo prende abbia meno oggetti totali.} \\
\\ 
\texttt{Sono d'accordo con questa divisione: tu prendi 2 specchi e 3 coltelli, io prendo 1 mattone e 1 pettine.} \\
\\ 
\texttt{Ti va bene se procediamo con questa proposta?} \\
            }
        }
    }
    & \\ \\

    \theutterance \stepcounter{utterance}  
    & & & \multicolumn{4}{p{0.6\linewidth}}{
        \cellcolor[rgb]{0.9,0.9,0.9}{
            \makecell[{{p{\linewidth}}}]{
                \texttt{\tiny{[GM$\langle$P2]}}
                \texttt{Perfetto! Sono contento che siamo riusciti a trovare una divisione che funziona bene per entrambi. } \\
\\ 
\texttt{Quindi per riassumere:} \\
\texttt{{-} Io: 2 specchi e 3 coltelli} \\
\texttt{{-} Tu: 1 mattone e 1 pettine} \\
\\ 
\texttt{Mi sembra una soluzione equilibrata che tiene conto delle nostre preferenze. Procediamo pure con questa divisione!} \\
\\ 
\texttt{{[}Proposta: 2 specchi, 3 coltelli{]}} \\
            }
        }
    }
     \\ \\

    \theutterance \stepcounter{utterance}  
    & & & \multicolumn{2}{p{0.3\linewidth}}{
        \cellcolor[rgb]{0.9,0.9,0.9}{
            \makecell[{{p{\linewidth}}}]{
                \texttt{\tiny{[GM$|$GM]}}
                \texttt{proposal} \\
            }
        }
    }
    & & \\ \\

    \theutterance \stepcounter{utterance}  
    & & & \multicolumn{2}{p{0.3\linewidth}}{
        \cellcolor[rgb]{0.9,0.9,0.9}{
            \makecell[{{p{\linewidth}}}]{
                \texttt{\tiny{[GM$|$GM]}}
                [2, 0, 3, 0]
            }
        }
    }
    & & \\ \\

    \theutterance \stepcounter{utterance}  
    & & \multicolumn{4}{p{0.6\linewidth}}{
        \cellcolor[rgb]{0.9,0.9,0.9}{
            \makecell[{{p{\linewidth}}}]{
                \texttt{\tiny{[P1$\langle$GM]}}
                \texttt{L'altro giocatore ha presentato una proposta segreta. Ora tocca a Lei presentare una sua proposta.} \\
            }
        }
    }
    & \\ \\

    \theutterance \stepcounter{utterance}  
    & \multicolumn{4}{p{0.6\linewidth}}{
        \cellcolor[rgb]{0.9,0.9,0.9}{
            \makecell[{{p{\linewidth}}}]{
                \texttt{\tiny{[P1$\rangle$GM]}}
                \texttt{{[}Proposta: 1 mattone, 1 pettine{]}} \\
            }
        }
    }
    & & \\ \\

    \theutterance \stepcounter{utterance}  
    & & & \multicolumn{2}{p{0.3\linewidth}}{
        \cellcolor[rgb]{0.9,0.9,0.9}{
            \makecell[{{p{\linewidth}}}]{
                \texttt{\tiny{[GM$|$GM]}}
                \texttt{proposal} \\
            }
        }
    }
    & & \\ \\

    \theutterance \stepcounter{utterance}  
    & & & \multicolumn{2}{p{0.3\linewidth}}{
        \cellcolor[rgb]{0.9,0.9,0.9}{
            \makecell[{{p{\linewidth}}}]{
                \texttt{\tiny{[GM$|$GM]}}
                [0, 1, 0, 1]
            }
        }
    }
    & & \\ \\

    \theutterance \stepcounter{utterance}  
    & & & \multicolumn{2}{p{0.3\linewidth}}{
        \cellcolor[rgb]{0.9,0.9,0.9}{
            \makecell[{{p{\linewidth}}}]{
                \texttt{\tiny{[GM$|$GM]}}
                [[0, 1, 0, 1], [2, 0, 3, 0]]
            }
        }
    }
    & & \\ \\

\end{supertabular}
}

\end{document}
